\section{Bayesian neural networks (Bayes by Backprop)}
\label{sec:bayes-by-backprop}

Section~\ref{subsec:variational-inference} introduced variational inference as a
general strategy rather than as a particular algorithm. Bayes by Backprop (BBB) is
one concrete instance of that strategy, applied to the weights of a neural
network~\cite{blundell2015}. The contrast with Section~\ref{sec:gaussian-processes}
is based on the assumption that the Gaussian process performs inference in the function space and obtains the predictive distribution in closed form, whereas the BBB performs inference in
the parameter space and settles for an approximation. The following therefore describes a
realization of the mathematical foundations of Chapter~\ref{ch:methods}, and the derivations of
Section~\ref{subsec:variational-inference} are not repeated.

Each weight of the network is replaced by a pair of variational parameters instead
of a single number $w_j$, the model stores a mean $\mu_j$ and a scale parameter
$\rho_j$, from which the standard deviation is obtained as   $\sigma_j = \log\!\left(1 + \exp \rho_j\right)$. The softplus transform keeps $\sigma_j$ positive without a constraint on the optimizer. The resulting variational family is exactly the mean-field Gaussian of Equation~\eqref{eq:mean-field}, with $\theta = (\mu, \rho)$ collecting all variational parameters. Under this parameterization, training a single network is replaced by training an infinite
ensemble whose members draw their weights from one shared, learnt distribution~\cite{blundell2015}.

% [Miejsce na Rysunek 3.4: waga deterministyczna kontra waga jako rozkład.
%  Dwa panele obok siebie, ten sam mały graf sieci (2-3-1). Po lewej każda
%  krawędź opisana jedną liczbą, po prawej — małą krzywą gaussowską o środku
%  mu_j i szerokości sigma_j. Cel rysunku: pokazać bez wzoru, na czym polega
%  różnica między punktowym oszacowaniem a rozkładem nad wagami.
%  UWAGA: rysunek narysować SAMODZIELNIE. Odpowiedniki u Blundella (rys. 1)
%  i Jospina (rys. 3) są objęte prawem autorskim i nie wolno ich przedrukować.]

\begin{figure}[htbp]
    \centering
    \includegraphics[width=0.85\textwidth]{rodzial3_rys/img3_4.png}
    \caption{Weight representation in a deterministic network and in a Bayesian
    neural network trained with Bayes by Backprop. On the left, each connection
    carries a single value obtained by point estimation. On the right, each
    connection carries a Gaussian distribution described by a mean $\mu_j$ and a
    standard deviation $\sigma_j$, so the number of stored parameters doubles.}
    \label{fig:bbb-weight-distributions}
\end{figure}

The objective function follows directly from Equation~\eqref{eq:elbo}, rewritten in
the notation used for network weights. Maximizing the evidence lower bound is
equivalent to minimizing the variational free energy,
%.
\begin{equation}
    \mathcal{F}(\mathcal{D}, \theta)
    = \underbrace{\mathrm{KL}\bigl(q(\mathbf{w} \mid \theta) \,\|\, p(\mathbf{w})\bigr)}_{\text{complexity cost}}
      - \underbrace{\mathbb{E}_{q(\mathbf{w} \mid \theta)}
        \bigl[\log p(\mathcal{D} \mid \mathbf{w})\bigr]}_{\text{likelihood cost}} ,
    \label{eq:bbb-free-energy}
\end{equation}
%
where $\mathbf{w}$ denotes a sample of all network weights~\cite{blundell2015}. The
likelihood cost rewards are in agreement with the training data. The complexity cost
penalizes departure from the prior and admits an interpretation as a compression
cost in the sense of minimum description length.

The second term of Equation (complexity cost) ~\eqref{eq:bbb-free-energy} deserves emphasis, because it settles a question raised in Section~\ref{subsec:mle-map}. Weight decay was shown
there to be MAP estimation under a Gaussian prior, that is, a regularization term
constructed on fact from a Bayesian reading. Here, the relationship is reversed - the penalty is not added to the loss by design but follows from the variational objective itself. Regularisation in a Bayesian neural network is therefore derived rather than imposed, and the strength of the penalty is fixed by the choice of prior rather than by a tuned coefficient~\cite{blundell2015}.

% ---------------------------------------------------------------------
\subsection{Training with unbiased Monte Carlo gradients}
\label{subsec:bbb-training}

The obstacle to optimizing Equation~\eqref{eq:bbb-free-energy} is not the objective, but the gradient. The weight sampling introduces stochastic nodes inside the network, and the backpropagation cannot pass through them in the usual way~\cite{jospin2022}. The expectation in the likelihood cost is taken over a distribution that itself depends on the parameters being optimised, so the gradient  cannot be moved inside the expectation without further argument. The
reparameterisation approach of Equation~\eqref{eq:reparameterization} removes the difficulty by writing a weight sample as a deterministic function of the variational parameters and of parameter-free noise~\cite{kingma2014}. 
For the Gaussian family adopted here, this takes the form
%
\begin{equation}
    \mathbf{w} = \mu + \log\!\left(1 + \exp \rho\right) \odot \varepsilon ,
    \qquad \varepsilon \sim \mathcal{N}(0, I) ,
    \label{eq:bbb-reparameterisation}
\end{equation}
%
where $\odot$ denotes multiplication in the element (elementwise). Randomness is limited to
$\varepsilon$, and gradients with respect to $\mu$ and $\rho$ pass through ordinary
backpropagation.

Variational treatments of neural networks predate this construction. Blundell et
al.\ build on the work of Graves (2011) and show that the gradients used there can
be made unbiased, and that the method can be extended to non-Gaussian priors~\cite{blundell2015}. Both improvements evolved from a single decision: the complexity cost is not evaluated in closed form. The objective is instead estimated entirely by sampling,
%
\begin{equation}
    \mathcal{F}(\mathcal{D}, \theta)
    \approx \sum_{i=1}^{n}
      \left[
        \log q\bigl(\mathbf{w}^{(i)} \mid \theta\bigr)
        - \log p\bigl(\mathbf{w}^{(i)}\bigr)
        - \log p\bigl(\mathcal{D} \mid \mathbf{w}^{(i)}\bigr)
      \right] ,
    \label{eq:bbb-mc-cost}
\end{equation}
%
where $\mathbf{w}^{(i)}$ is the $i$-th sample drawn from the variational posterior.
Dispensing with the closed form frees the choice of prior. Blundell et al.\ exploit
this freedom with a scale mixture of two zero-mean Gaussians, the first broad and
the second sharply concentrated around zero. On a network with two hidden layers of
1200 units, this prior reached a test error of 1.32\% on MNIST, against 2.04\% for a
plain Gaussian prior.

Equation~\eqref{eq:bbb-mc-cost} has another less obvious property - every term is
evaluated on the same weight sample, a variance reduction technique known as common
random numbers~\cite{blundell2015}. Where part of the objective is available in
closed form, that part is insensitive to the particular draw, while the remainder is
not. A qualification is in order, and it comes from the authors themselves. The
sampled complexity cost was not found to outperform a closed-form Kullback--Leibler
term where the latter could be computed, only to perform no worse. The argument for
the sampling estimator rests on the range of admissible priors, not on measured
accuracy.

The resulting training loop remains close to that of a deterministic
network~\cite{jospin2022}. Each step draws $\varepsilon$, assembles the weights
through Equation~\eqref{eq:bbb-reparameterisation}, evaluates the objective of
Equation~\eqref{eq:bbb-mc-cost} on a minibatch, and updates $\mu$ and $\rho$ by
gradient descent. The gradient with respect to the sampled weights is shared between
both updates and is precisely what ordinary backpropagation returns~\cite{blundell2015}.
Standard optimizers such as Adam apply without modification. One further detail
concerns minibatches: the complexity cost refers to the whole data set and must be
distributed between them, for instance, by weighting it with coefficients
$\pi_i \in [0,1]$ satisfying $\sum_{i=1}^{M} \pi_i = 1$,
%
\begin{equation}
    \mathcal{F}_i^{\pi}(\mathcal{D}_i, \theta)
    = \pi_i\, \mathrm{KL}\bigl(q(\mathbf{w} \mid \theta) \,\|\, p(\mathbf{w})\bigr)
      - \mathbb{E}_{q(\mathbf{w} \mid \theta)}
        \bigl[\log p(\mathcal{D}_i \mid \mathbf{w})\bigr] .
    \label{eq:bbb-kl-reweighting}
\end{equation}
%
A schedule that front-loads the complexity cost onto the first minibatches of each
epoch was reported to work well~\cite{blundell2015}. It should be noted that the
objective is estimated from a small number of samples, often a single one, so the
gradient estimate is noisy and the convergence curve is far less smooth than in
deterministic training~\cite{jospin2022}. Averaging the reported loss over several
epochs gives a more readable approach of convergence.

% ---------------------------------------------------------------------
\subsection{Prediction and uncertainty decomposition}
\label{subsec:bbb-prediction}

The prediction reproduces the marginalization of Equation~\eqref{eq:predictive-2} by
Monte Carlo integration. A test input is passed through the network $T$ times, and
each pass draws an independent weight sample $\mathbf{w}^{(t)}$ from the variational
posterior. 
Each pass yields a predictive mean $\mu_{\mathbf{w}^{(t)}}(x^\ast)$ and a
predictive variance $\sigma^2_{\mathbf{w}^{(t)}}$, which under the homoscedastic
noise model adopted here does not depend on the input. The prediction reported for $x^\ast$ is the
average across passes~\cite{jospin2022},
%
\begin{equation}
    \hat{\mu}(x^\ast)
    = \frac{1}{T} \sum_{t=1}^{T} \mu_{\mathbf{w}^{(t)}}(x^\ast) .
    \label{eq:bbb-predictive-mean}
\end{equation}

The spread across passes is not a by-product of this averaging but the quantity of
primary interest. Both components of Equation~\eqref{eq:total-variance-2} are
estimated from the same $T$ samples,
%
\begin{equation}
    \operatorname{Var}\!\left(y^\ast \mid x^\ast\right)
    \approx
    \underbrace{\frac{1}{T} \sum_{t=1}^{T}
      \sigma^2_{\mathbf{w}^{(t)}}(x^\ast)}_{\text{aleatoric}}
    +
    \underbrace{\frac{1}{T-1} \sum_{t=1}^{T}
      \bigl(\mu_{\mathbf{w}^{(t)}}(x^\ast) - \hat{\mu}(x^\ast)\bigr)^2}_{\text{epistemic}} .
    \label{eq:bbb-variance-estimator}
\end{equation}
%
The first term averages the noise predicted by the individual networks, while the second
measures how far those networks disagree about the mean, and it vanishes when the
variational posterior collapses to a point. Where the training data are dense, the
sampled networks converges and the epistemic term is small, but if there is absent data, the
weight distributions are constrained only by the prior, the sampled networks diverge, and the epistemic term grows.

% [Miejsce na Rysunek 3.5: posterior BBB na zbiorze load_sin.
%  Trening na [0, 6], ewaluacja na [-2, 10]. TE SAME osie i ta sama skala co
%  Rysunek 3.2 (posterior GP), żeby oba dały się porównać wzrokiem.
%  Predykcyjna średnia z Równania (bbb-predictive-mean), pasmo +-2 sigma
%  z pierwiastka Równania (bbb-variance-estimator), punkty treningowe,
%  prawdziwa funkcja linią przerywaną. Opcjonalnie: 10-15 pojedynczych
%  przebiegów cienką linią, żeby widać było rozrzut, z którego liczony jest
%  człon epistemiczny. Cel rysunku: czy pasmo rozchodzi się poza zakresem
%  treningowym słabiej niż u GP.]

\begin{figure}[htbp]
    \centering
    \includegraphics[width=0.85\textwidth]{rodzial3_rys/img3_5.png}
    \caption{Predictive distribution of a Bayesian neural network trained with Bayes
    by Backprop on the synthetic sine dataset. The model is trained on the interval
    $[0, 6]$ and evaluated on $[-2, 10]$, with the axes matching those of
    Figure~\ref{fig:gp-posterior-sine} to allow a direct comparison. The band shows
    $\pm 2\sigma$ obtained from Equation~\eqref{eq:bbb-variance-estimator} over
    $T$ stochastic forward passes.}
    \label{fig:bbb-posterior-sine}
\end{figure}

% ---------------------------------------------------------------------
\subsection{Computational cost and limitations}
\label{subsec:bbb-limitations}
 
The cost of the method is incurred twice. Training stores two variational parameters
per weight, so the parameter count doubles relative to a deterministic network of the
same architecture~\cite{arbel2023}. Each training iteration was reported to run about
twice as slowly as an iteration with dropout~\cite{blundell2015}. Inference is
affected more severely because a single forward pass no longer suffices and the
$T$ passes of Equation~\eqref{eq:bbb-variance-estimator} must be performed for every
test point. The need for Monte Carlo integration at evaluation time is a persistent
drawback of Bayesian neural networks and not a property of this algorithm
alone~\cite{jospin2022}.

The mean-field assumption is the most consequential limitation. Learning a complete
covariance matrix on the weights would require $\mathcal{O}(n^2)$ variational
parameters, which is unattainable for the parameter counts reported in
Section~\ref{subsec:intractability}; the diagonal approximation reduces this to
$\mathcal{O}(n)$~\cite{jospin2022}. The independence between weights is therefore a
concession to tractability rather than a modeling belief. Its cost is stated
plainly in the literature: the restriction is likely to result in an underestimate
of the overall posterior uncertainty~\cite{arbel2023}. Variational methods are also
prone to mode collapse, that is, to concentrating on a single mode of a posterior
known to be highly multimodal. Both effects act in the same direction and narrow the
approximate posterior.

A qualification is again due, and again it comes from the original work. Blundell et
al.\ acknowledge the known tendency of variational methods to underestimate
uncertainty, yet report that they did not observe its practical consequences in
their experiments~\cite{blundell2015}. The claim is therefore a prediction rather
than a settled result. Equation~\eqref{eq:total-variance-2} identifies where such an
underestimate would appear: the epistemic term of Equation~\eqref{eq:bbb-variance-estimator} is the only component affected by the shape of the approximate posterior. This gives a testable expectation for Chapter~\ref{ch:results}. If the mean-field posterior is too narrow, Bayes by
Backprop should produce narrower prediction intervals than the Gaussian process of
Section~\ref{sec:gaussian-processes}, and its empirical coverage should fall short of
the nominal level. The metrics defined in Chapter~\ref{ch:setup} are reported for
both methods precisely to settle this question.
